\section{Выводы}

Выполнив данную лабораторную работу по курсу \enquote{Дискретный анализ},
я реализовал алгоритм Кнута"---Морриса"---Пратта для поиска образца в тексте над
алфавитом целых чисел. Разобрался с построением префикс-функции и механизмом
откатов, позволяющим избежать повторных сравнений. Реализовал корректную обработку
ведущих нулей через стандартное считывание \texttt{uint32\_t}, а также потоковый
поиск вхождений, пересекающих границы строк, без хранения всего текста в памяти.

Тестирование подтвердило теоретические оценки: на случайных данных преимущество КМП
не всегда заметно из-за коротких частичных совпадений, однако на патологических входах
(образец из 19 единиц и завершающей двойки в тексте из единиц) алгоритм работает
примерно в \textbf{4.7$\times$} быстрее наивного поиска при $n = 5\,000\,000$,
демонстрируя линейное время против квадратичного.

\pagebreak
